% !PRINTED 420
\weberpage{420}
Now in \S.\,56 we have become acquainted with two forms $f$ and $H$ of degrees $6$ and $8$, which remain absolutely unchanged by the tetrahedral substitutions when these are taken homogeneous and with determinant 1. To be sure, we cannot apply these forms directly in the form in which they are given there, because there the tetrahedral group was assumed in a different shape. But we find the first of these forms very simply if we recall that its roots were the two-fold poles of the tetrahedral group, and hence the roots of the three equations
\[
x = \psi(x), \quad x = \chi(x), \quad x = \psi\chi(x)
\]
or
\[
x^2 + 1 = 0, \quad x^2 - 2\omega x - 1 = 0, \quad \omega x^2 + 2x - \omega = 0,
\]
and thus, with the use of the relation
\[
\omega^2 + \omega = 1
\]
one obtains the equation for the two-fold poles
\[
(x^2+1)(x^4+2x^3-6x^2-2x+1) = 0.
\]

Hence the sought tetrahedral form of degree $6$ in the homogeneous variables $x_1, x_2$ becomes
\weq{5}{t(x_1,x_2) = x_1^6 + 2x_1^5x_2 - 5x_1^4x_2^2 - 5x_1^2x_2^4 - 2x_1x_2^5 + x_2^6.}

The second tetrahedral form that we seek is the \gesperrt{Hesse} determinant of this.

If we set
\[
\tilde\omega(x_1,x_2) = \dfrac{1}{5^2 \cdot 4^2} \left\{ t''(x_1,x_1)\, t''(x_2,x_2) - [t''(x_1,x_2)]^2 \right\},
\]
then a simple calculation yields
\weq{6}{\tilde\omega(x_1,x_2) = -(x_1^8+x_2^8) + (x_1^7x_2 - x_1x_2^7) - 7(x_1^6x_2^2+x_1^2x_2^6) - 7(x_1^5x_2^3-x_1^3x_2^5).}

We further adduce the three icosahedral forms of degrees $12$, $20$ and $30$, which indeed also remain unchanged by the tetrahedral substitutions (\S.\,60):
\weq{7}{f(x_1,x_2) = x_1x_2(x_1^{10}+11x_1^5x_2^5-x_2^{10})}
\weq{8}{H(x_1,x_2) = -(x_1^{20}+x_2^{20}) + 228(x_1^{15}x_2^5-x_1^5x_2^{15}) - 494x_1^{10}x_2^{10}}
\weq{9}{T(x_1,x_2) = (x_1^{30}+x_2^{30}) + 522(x_1^{25}x_2^5-x_1^5x_2^{25}) - 10005(x_1^{20}x_2^{10}+x_2^{20}x_1^{10}).}
